% CS 4023/5023 Project 1 -- Robot Code Documentation (Section 5.4)
% OUTLINE: each section lists what to write. Replace the bullet prompts
% with prose. Target length 1.5--2 pages of text (about 80 characters per
% line, 50 lines per page); figures and tables do not count toward it.

\documentclass[11pt]{article}
\usepackage[margin=1in]{geometry}
\usepackage{booktabs}
\usepackage{graphicx}
\usepackage{hyperref}

\newcommand{\prompt}[1]{\textit{[#1]}}

\title{Project 1: Robot Code Documentation}
\author{Team members: \prompt{names}}
\date{\prompt{date}}

\begin{document}
\maketitle

% ---------------------------------------------------------------------
\section{Overview}
\begin{itemize}
  \item \prompt{One paragraph: what the code does.} Two ROS~2 Python nodes in
        package \texttt{proj1}: \texttt{reactive\_node.py}
        (class \texttt{SubsumptionController}) drives the TurtleBot~4;
        \texttt{occupancy\_mapper.py} (class \texttt{OccupancyMapper})
        builds the map. They run independently and share only sensor topics.
  \item \prompt{Figure 1: node/topic graph} --
        \texttt{/scan}, \texttt{/odom}, \texttt{/hazard\_detection},
        \texttt{/teleop\_cmd} $\rightarrow$ controller $\rightarrow$
        \texttt{/cmd\_vel}; \texttt{/scan}, \texttt{/odom}, TF
        $\rightarrow$ mapper $\rightarrow$ \texttt{/map}.
\end{itemize}

% ---------------------------------------------------------------------
\section{Reactive Architecture}
\subsection{Which architecture}
\begin{itemize}
  \item \prompt{Name it.} Priority-based pseudo-subsumption in the style of
        Jones et al.: every behavior proposes a command or abstains, and the
        highest-priority proposal wins.
  \item \prompt{Contrast briefly with Brooks' subsumption (suppression and
        inhibition wires between layers) and with schema theory (vector
        summation of motor schemas), and say what we did not do.}
\end{itemize}

\subsection{Why we chose it}
\begin{itemize}
  \item \prompt{Reasons, e.g.: the brief gives a strict priority ordering,
        which maps directly onto a fixed arbitration list; behaviors are
        independent and easy to test one at a time; no vector tuning as in
        schemas; simple to explain and debug.}
\end{itemize}

\subsection{Where the code embodies it}
\begin{itemize}
  \item \texttt{self.behaviors} tuple: the priority order, highest first.
  \item \texttt{control\_loop()}: the arbiter; publishes the first
        non-\texttt{None} command each cycle (10~Hz).
  \item Each behavior is one method returning a command or \texttt{None}.
  \item \prompt{Figure 2: priority stack diagram, halt on top,
        forward at the bottom, arbiter on the right.}
\end{itemize}

% ---------------------------------------------------------------------
\section{Behaviors}
\prompt{One short paragraph per behavior, or rely mostly on the table.
Point out reflex vs.\ fixed action pattern explicitly.}

\begin{table}[h]
\centering\small
\begin{tabular}{@{}llll@{}}
\toprule
\# & Method & Trigger & Response / type \\
\midrule
1 & \texttt{halt} & any bump in \texttt{/hazard\_detection} & stop; reflex \\
2 & \texttt{teleop} & key command $<0.5$\,s old & pass through \\
3 & \texttt{escape} & both front arcs $<1$\,ft, nearly equal & turn
    $180^\circ\pm30^\circ$; FAP \\
4 & \texttt{avoid} & either front arc $<1$\,ft & turn away from closer
    side; reflex \\
5 & \texttt{random\_turn} & 1\,ft of forward travel & turn
    $U(-15^\circ,15^\circ)$ \\
6 & \texttt{forward} & always & drive at 0.2\,m/s \\
\bottomrule
\end{tabular}
\caption{Behaviors, highest priority first.}
\end{table}

\begin{itemize}
  \item \prompt{Explain why escape is a fixed action pattern: the target
        heading is latched in \texttt{turn\_target}/\texttt{turn\_owner}
        and the turn finishes even after the obstacle disappears.}
  \item \prompt{Explain why avoid is a reflex: it is recomputed every
        cycle and stops as soon as nothing is within 1~ft.}
  \item \prompt{Note that avoid cancels a pending random turn.}
\end{itemize}

% ---------------------------------------------------------------------
\section{Data Structures}
\subsection{Controller}
\begin{itemize}
  \item Sensor state: \texttt{bumper\_hit}, \texttt{front\_left\_dist},
        \texttt{front\_right\_dist}, cached \texttt{scan\_yaw\_offset}.
  \item Odometry state: \texttt{current\_yaw}, last position,
        \texttt{forward\_travel} accumulator.
  \item Teleop state: last twist and its arrival time.
  \item Latched turn: \texttt{turn\_target} (absolute heading) and
        \texttt{turn\_owner} (which behavior owns it).
  \item ROS parameters for every threshold (speeds, 1~ft distance,
        $\pm30^\circ$ front arc, symmetry tolerance, turn angles).
        \prompt{Say why: the TA can retune without editing code.}
\end{itemize}

\subsection{Mapper}
\begin{itemize}
  \item 2-D NumPy \texttt{float32} array of log-odds, 0.05\,m cells,
        20\,m $\times$ 20\,m centered on the start pose.
  \item Boolean \texttt{seen} mask so never-observed cells stay unknown
        ($-1$).
  \item Cached lidar mount pose $(x, y, \text{yaw})$ in the base frame.
\end{itemize}

% ---------------------------------------------------------------------
\section{Algorithms}
\begin{itemize}
  \item \textbf{Front-obstacle detection:} convert each beam angle into the
        robot frame using the lidar's TF yaw offset, keep the minimum range
        in the left ($0$ to $+30^\circ$) and right ($-30^\circ$ to $0$)
        arcs. \prompt{Why TF: robust if the TA moves or rotates the lidar.}
  \item \textbf{Symmetry test:} both arcs $<$ 1~ft and
        $|d_L - d_R| <$ tolerance $\Rightarrow$ escape; otherwise avoid.
  \item \textbf{Distance for random turn:} project each odometry step onto
        the heading; only positive (forward) motion counts.
  \item \textbf{Latched turns:} target = current yaw + angle, wrapped to
        $[-\pi,\pi]$; rotate at fixed speed until within $3^\circ$.
  \item \textbf{Occupancy mapping:} per scan, transform the lidar pose into
        the odom frame, ray-cast each beam at half-cell steps, add
        $L_{\text{free}}$ to traversed cells and $L_{\text{occ}}$ to the end
        cell, clamp to $\pm5$, publish $p = 1 - 1/(1+e^{l})$ scaled to
        0--100. \prompt{Cite the log-odds occupancy-grid formulation.}
\end{itemize}

% ---------------------------------------------------------------------
\section{Design Decisions and Limitations}
\begin{itemize}
  \item \prompt{Map is in the odom frame, so it drifts with odometry; no
        localization, as the brief does not require navigation.}
  \item \prompt{``1~ft in front'' is measured from the lidar.}
  \item \prompt{Halt outranks teleop, so a pressed bumper blocks keyboard
        control too.}
  \item \prompt{Create~3 built-in bump reflexes are disabled at launch so
        that halt is in charge.}
  \item \prompt{Coding standards: ROS~2 Python style, flake8 and pep257
        tests, object-oriented nodes.}
\end{itemize}

% ---------------------------------------------------------------------
\section{Sources and Credit}
\prompt{Required by Section~6 of the assignment. List every source of code
or text and which parts are original, changed, or taken from where
(including tools used to help write code or documents). List who on the
team wrote which parts.}

\end{document}
